% ============================================================
% Appendice D — Risorse e letture consigliate
% ============================================================

\chapter{Risorse e Letture Consigliate}
\label{app:risorse}

Riferimenti per approfondire i temi trattati nel libro.

\begin{notebox}
I link a documentazione ufficiale di agenti AI cambiano frequentemente. Verifica che gli URL siano ancora validi prima di seguirli.
\end{notebox}

\section*{Il Framework ADLC}

\begin{itemize}
  \item \textbf{Repository ufficiale:} \url{https://github.com/rollaradio/adlc-framework} --- sorgente del framework, CHANGELOG, MIGRATION.md.
  \item \textbf{MANUAL.it.md} (nel repository): manuale tecnico compatto per riferimento rapido.
  \item \textbf{COMMANDS.md}: specifica formale di tutti i comandi conversazionali.
\end{itemize}

\section*{Prompt Engineering e Contesto}

\begin{itemize}
  \item \textbf{The Prompt Report} --- Schulhoff et al. (2024). Survey sulle tecniche di prompt engineering. Spiega perché il contesto strutturato funziona meglio delle istruzioni inline.
  \item \textbf{Chain-of-Thought Prompting Elicits Reasoning in LLMs} --- Wei et al. (2022). Spiega perché ``pianificare prima di agire'' produce output migliori.
  \item \textbf{Software Design for AI (SDD)} --- Fowler, martinfowler.com. Prospettiva sull'integrazione degli agenti AI nei processi di sviluppo.
\end{itemize}

\section*{Sicurezza}

\begin{itemize}
  \item \textbf{OWASP Top 10 (2021):} \url{https://owasp.org/www-project-top-ten/}. I vincoli SEC-01, SEC-02, SEC-08, SEC-09 derivano direttamente da OWASP.
  \item \textbf{OWASP ASVS:} \url{https://owasp.org/www-project-application-security-verification-standard/}. Standard dettagliato per vincoli SEC avanzati.
  \item \textbf{OWASP Logging Cheat Sheet}: riferimento diretto per SEC-08.
  \item \textbf{OWASP SSRF Prevention Cheat Sheet}: riferimento per SEC-09.
\end{itemize}

\section*{Testing e Qualità}

\begin{itemize}
  \item \textbf{The Testing Pyramid} --- Martin Fowler, \url{https://martinfowler.com/articles/practical-test-pyramid.html}. Guida la strategia di test del libro.
  \item \textbf{EPAM Testing Pyramid for AI-Assisted Development} --- EPAM Systems (2024). Adattamento della piramide per contesti con agenti AI.
\end{itemize}

\section*{AI-Assisted Development}

\begin{itemize}
  \item \textbf{GitHub Copilot Impact Study} --- GitHub/Accenture (2024). Ricerca quantitativa sull'impatto di Copilot sulla produttività.
  \item \textbf{Salesforce AI Developer Survey} --- Salesforce (2024). Evidenzia i problemi di context drift e vincoli dimenticati.
  \item \textbf{Cycode State of AI-Assisted Development} --- Cycode (2024). Analisi dei rischi di sicurezza con agenti AI non strutturati.
\end{itemize}

\section*{Architettura}

\begin{itemize}
  \item \textbf{Architectural Decision Records} --- Michael Nygard: \url{https://cognitect.com/blog/2011/11/15/documenting-architecture-decisions}. Il formato ADR usato per decisioni CRITICAL.
\end{itemize}

\section*{Standard di riferimento}

\begin{itemize}
  \item \textbf{NIST SP 800-53}: standard di controlli di sicurezza per sistemi informativi. I vincoli SEC di ADLC ne sono parzialmente ispirati.
  \item \textbf{Google SRE Book:} \url{https://sre.google/sre-book/table-of-contents/}. Capitoli su SLO/SLI (base di PERF-01/02) e incident response (base di \texttt{06\_OPS.md}).
\end{itemize}

\section*{Documentazione agenti AI}

\begin{itemize}
  \item \textbf{Claude Code --- Anthropic:} \url{https://docs.anthropic.com/claude-code}
  \item \textbf{GitHub Copilot:} \url{https://docs.github.com/en/copilot}
  \item \textbf{OpenAI CLI / Codex} (verifica stato attuale --- il prodotto è soggetto a variazioni): \url{https://github.com/openai/openai-codex}
\end{itemize}
